\documentclass[10pt,a4paper]{amsart}
\usepackage{amsmath,amssymb,amsthm,graphicx}
% Retained preamble for the Stacks Project chapter snapshot in this directory.
% The chapter's own preamble.tex is a copy of the shared Stacks Project preamble; the block
% below is sliced MECHANICALLY, line for line, from it. Nothing was added, removed or
% redefined; the wrapper class, the xr-hyper/hyperref wiring and the externaldocument list
% are omitted because the assembler supplies the document class.
%
% lines taken from preamble.tex: 3, 5, 9, 11, 12, 156, 157, 158, ... (51 lines total)
%
% package coverage restored from the chapter's own preamble (2026-10-13): xy, xy-2cell
% package coverage restored from the chapter's own preamble (2026-10-14): verbatim
%   the chapter's preamble.tex loads \usepackage{verbatim} on line 8 and then defines
%   reference/slogan/history as \comment environments; without the package \comment is
%   undefined and any card whose counted spans use one of the three environments fails
%   with "Undefined control sequence" (L44184, lemma-representable-over-space).

\usepackage{verbatim}
\usepackage[all]{xy}
\xyoption{2cell}
\UseAllTwocells

\newenvironment{reference}{\comment}{\endcomment}
\newenvironment{slogan}{\comment}{\endcomment}
\newenvironment{history}{\comment}{\endcomment}
\theoremstyle{plain}
\newtheorem{theorem}[subsection]{Theorem}
\newtheorem{proposition}[subsection]{Proposition}
\newtheorem{lemma}[subsection]{Lemma}
\theoremstyle{definition}
\newtheorem{definition}[subsection]{Definition}
\newtheorem{example}[subsection]{Example}
\newtheorem{exercise}[subsection]{Exercise}
\newtheorem{situation}[subsection]{Situation}
\theoremstyle{remark}
\newtheorem{remark}[subsection]{Remark}
\newtheorem{remarks}[subsection]{Remarks}
\numberwithin{equation}{subsection}
\def\lim{\mathop{\mathrm{lim}}\nolimits}
\def\colim{\mathop{\mathrm{colim}}\nolimits}
\def\Spec{\mathop{\mathrm{Spec}}}
\def\Hom{\mathop{\mathrm{Hom}}\nolimits}
\def\Ext{\mathop{\mathrm{Ext}}\nolimits}
\def\SheafHom{\mathop{\mathcal{H}\!\mathit{om}}\nolimits}
\def\SheafExt{\mathop{\mathcal{E}\!\mathit{xt}}\nolimits}
\def\Sch{\mathit{Sch}}
\def\Mor{\mathop{\mathrm{Mor}}\nolimits}
\def\Ob{\mathop{\mathrm{Ob}}\nolimits}
\def\Sh{\mathop{\mathit{Sh}}\nolimits}
\def\NL{\mathop{N\!L}\nolimits}
\def\CH{\mathop{\mathrm{CH}}\nolimits}
\def\proetale{{pro\text{-}\acute{e}tale}}
\def\etale{{\acute{e}tale}}
\def\QCoh{\mathit{QCoh}}
\def\Ker{\mathop{\mathrm{Ker}}}
\def\Im{\mathop{\mathrm{Im}}}
\def\Coker{\mathop{\mathrm{Coker}}}
\def\Coim{\mathop{\mathrm{Coim}}}
\DeclareMathSymbol{\boxtimes}{\mathbin}{AMSa}{"02}
\def\QCohstack{\mathcal{QC}\!\mathit{oh}}
\def\Cohstack{\mathcal{C}\!\mathit{oh}}
\def\Spacesstack{\mathcal{S}\!\mathit{paces}}
\def\Quotfunctor{\mathrm{Quot}}
\def\Hilbfunctor{\mathrm{Hilb}}
\def\Curvesstack{\mathcal{C}\!\mathit{urves}}
\usepackage{cleveref}
\crefname{originallemma}{original lemma}{original lemmas}
\Crefname{originallemma}{Original lemma}{Original lemmas}
\makeatletter
\newlabel{lemma-space-presentation}{{0262}{source}{}{sourcepointer-1}{}}
\newlabel{lemma-space-presentation@cref}{{[originallemma][0][]0262}{[1][0][]source}{}{}{}}
\makeatother
\crefname{originaltheorem}{original theorem}{original theorems}
\Crefname{originaltheorem}{Original theorem}{Original theorems}
\makeatletter
\newlabel{theorem-presentation}{{02WW}{source}{}{sourcepointer-2}{}}
\newlabel{theorem-presentation@cref}{{[originaltheorem][0][]02WW}{[1][0][]source}{}{}{}}
\makeatother
\crefname{originallemma}{original lemma}{original lemmas}
\Crefname{originallemma}{Original lemma}{Original lemmas}
\makeatletter
\newlabel{topologies-lemma-change-alpha}{{022K}{source}{}{sourcepointer-3}{}}
\newlabel{topologies-lemma-change-alpha@cref}{{[originallemma][0][]022K}{[1][0][]source}{}{}{}}
\makeatother
\crefname{originalsection}{original section}{original sections}
\Crefname{originalsection}{Original section}{Original sections}
\makeatletter
\newlabel{topologies-section-change-alpha}{{022I}{source}{}{sourcepointer-4}{}}
\newlabel{topologies-section-change-alpha@cref}{{[originalsection][0][]022I}{[1][0][]source}{}{}{}}
\makeatother
\title{Restriction to a smaller big site preserves schemes and algebraic spaces}
\author{The Stacks Project Authors}
\date{stacks-project:spaces.tex:local-snapshot-20261012; GNU FDL 1.2 or later}
\begin{document}
\maketitle
\setcounter{subsection}{1}
\noindent\textit{Editorial note.} All delivered mathematical text is the author's own, copied from the pinned Stacks Project chapter file spaces.tex; no statement or proof step was rewritten, reordered or removed.\par\medskip
\noindent\textit{Editorial note.} The statement contains an align block displaying the two morphisms of topoi; it is part of the counted statement span and is preserved verbatim.\par\medskip
\noindent\textit{Editorial note.} The four references of the counted spans are recorded as non-load-bearing navigation pointers; the argument that consumes them stays inside the counted proof span.\par\medskip
\noindent\textit{Editorial note.} Preamble: sources/stacks-spaces/retained-preamble.tex, sliced mechanically from the chapter's own preamble.tex.\par\medskip
\noindent\textit{Allowed standard prerequisites (editorial).}\begin{itemize}
\item Description of the change-of-big-site morphisms of topoi. Topologies, Section topologies-section-change-alpha and Lemma topologies-lemma-change-alpha, used at proof lines 2776-2789 and 2805-2807. Supplies f\textasciicircum{}\{-1\} and the identification f\textasciicircum{}\{-1\}h\_X = h'\_X.
\item Presentation of an algebraic space. Lemma lemma-space-presentation (1387-1407, produced as L44183), used at proof lines 2810-2813. Writes F = h\_U/h\_R.
\item The presentation theorem for algebraic spaces. Theorem theorem-presentation, used at proof lines 2814-2815. Concludes that f\textasciicircum{}\{-1\}F is an algebraic space.
\end{itemize}\medskip
\ifdefined\hypertarget\hypertarget{sourcepointer-1}{}\fi
\noindent\textit{Source reference.} Original lemma 0262: focused-seeds/AGP/sources/stacks-snapshot-20261011/tags.txt, source lines 2835--2835. Navigation pointer into the Stacks Project; the standard fact it supplies is applied inside the counted proof span and the argument resting on it stays there.\par\medskip
\ifdefined\hypertarget\hypertarget{sourcepointer-2}{}\fi
\noindent\textit{Source reference.} Original theorem 02WW: focused-seeds/AGP/sources/stacks-snapshot-20261011/tags.txt, source lines 3819--3819. Navigation pointer into the Stacks Project; the standard fact it supplies is applied inside the counted proof span and the argument resting on it stays there.\par\medskip
\ifdefined\hypertarget\hypertarget{sourcepointer-3}{}\fi
\noindent\textit{Source reference.} Original lemma 022K: focused-seeds/AGP/sources/stacks-snapshot-20261011/tags.txt, source lines 2709--2709. Navigation pointer into the Stacks Project; the standard fact it supplies is applied inside the counted proof span and the argument resting on it stays there.\par\medskip
\ifdefined\hypertarget\hypertarget{sourcepointer-4}{}\fi
\noindent\textit{Source reference.} Original section 022I: focused-seeds/AGP/sources/stacks-snapshot-20261011/tags.txt, source lines 2707--2707. Navigation pointer into the Stacks Project; the standard fact it supplies is applied inside the counted proof span and the argument resting on it stays there.\par\medskip
\begin{lemma}
\label{lemma-change-big-site}
Suppose given big sites $\Sch_{fppf}$ and $\Sch'_{fppf}$.
Assume that $\Sch_{fppf}$ is contained in $\Sch'_{fppf}$,
see Topologies, Section \ref{topologies-section-change-alpha}.
Let $S$ be an object of $\Sch_{fppf}$. Let
\begin{align*}
g : \Sh((\Sch/S)_{fppf})
\longrightarrow
\Sh((\Sch'/S)_{fppf}), \\
f : \Sh((\Sch'/S)_{fppf})
\longrightarrow
\Sh((\Sch/S)_{fppf})
\end{align*}
be the morphisms of topoi of
Topologies, Lemma \ref{topologies-lemma-change-alpha}.
Let $F$ be a sheaf of sets on $(\Sch/S)_{fppf}$. Then
\begin{enumerate}
\item if $F$ is representable by a scheme
$X \in \Ob((\Sch/S)_{fppf})$ over $S$,
then $f^{-1}F$ is representable too, in fact it is representable by the
same scheme $X$, now viewed as an object of $(\Sch'/S)_{fppf}$, and
\item if $F$ is an algebraic space over $S$, then $f^{-1}F$ is an algebraic
space over $S$ also.
\end{enumerate}
\end{lemma}

\begin{proof}
Let $X \in \Ob((\Sch/S)_{fppf})$. Let us write $h_X$ for the
representable sheaf on $(\Sch/S)_{fppf}$ associated to $X$, and
$h'_X$ for the representable sheaf on $(\Sch'/S)_{fppf}$ associated to
$X$. By the description of $f^{-1}$ in
Topologies, Section \ref{topologies-section-change-alpha}
we see that $f^{-1}h_X = h'_X$. This proves (1).

\medskip\noindent
Next, suppose that $F$
is an algebraic space over $S$. By Lemma \ref{lemma-space-presentation}
this means that $F = h_U/h_R$ for some \'etale equivalence relation
$R \to U \times_S U$ in $(\Sch/S)_{fppf}$. Since $f^{-1}$ is an
exact functor we conclude that $f^{-1}F = h'_U/h'_R$. Hence
$f^{-1}F$ is an algebraic space over $S$ by Theorem \ref{theorem-presentation}.
\end{proof}

\end{document}
